% Part: second-order-logic
% Chapter: syntax-and-semantics
% Section: expressive-power

\documentclass[../../../include/open-logic-section]{subfiles}

\begin{document}

\olfileid[fr]{sol}{syn}{exp}
\olsection{Puissance expressive}

\begin{explain}
La quantification sur les variables du second ordre
accroît considérablement la puissance expressive
du langage par rapport à la logique du premier
ordre. La quantification existentielle du second
ordre permet d'affirmer qu'il existe des fonctions
ou des relations possédant certaines propriétés.
Au premier ordre, il faut, pour cela, prévoir un
symbole non logique, c'est-à-dire un
!!{function} ou un !!{predicate}. La quantification
universelle du second ordre permet d'affirmer que
tous les sous-ensembles du !!{domain}, toutes les
relations sur ce !!{domain}, ou toutes les fonctions
de ce !!{domain} dans lui-même possèdent une
propriété. Au premier ordre, on ne peut parler que
des sous-ensembles, relations ou fonctions attribués
à l'un des symboles non logiques du langage.
De plus, pour préciser la propriété dont on affirme
qu'elle appartient à certains ou à tous ces objets,
on peut encore employer la quantification du second
ordre. Cela permet de définir des relations et
d'exprimer des propriétés du !!{domain} qui ne sont
pas définissables ou exprimables au premier ordre.
\end{explain}

\begin{defn}
Soit $\Struct{M}$ une !!{structure} pour un
langage~$\Lang{L}$. Une relation
$R\subseteq\Domain{M}^2$ est
\emph{définissable} dans~$\Lang{L}$ s'il existe
une !!{formula}~$!A_R(x,y)$ dont les seules
!!{variable}s libres sont $x$ et~$y$, telle que,
pour tous $a,b\in\Domain{M}$ et toute
assignation~$s$ avec $s(x)=a$ et $s(y)=b$,
$R(a,b)$, c'est-à-dire $\tuple{a,b}\in R$,
si et seulement si $\Sat{M}{!A_R(x,y)}[s]$.
\end{defn}

\begin{ex}
Au premier ordre, la !!{formula}~$\eq[x][y]$
définit la relation d'identité
$\Id{\Domain{M}}$, c'est-à-dire
$\Setabs{\tuple{a,a}}{a\in\Domain{M}}$.
Au second ordre, nous pouvons la définir
\emph{sans~$\eq$}. Si $a$ et $b$ sont le même
!!{element} de~$\Domain{M}$, ils appartiennent
exactement aux mêmes sous-ensembles de ce
!!{domain}. Réciproquement, s'ils sont distincts,
ils n'appartiennent pas aux mêmes sous-ensembles :
par exemple, $a\in\{a\}$ mais $b\notin\{a\}$
si $a\neq b$. La propriété « appartenir aux
mêmes sous-ensembles de~$\Domain{M}$ » vaut donc
de $a$ et~$b$ exactement lorsque $a=b$.
Elle s'exprime au second ordre, puisqu'on peut
quantifier sur tous les sous-ensembles
de~$\Domain{M}$. La !!{formula} suivante définit
donc $\Id{\Domain{M}}$ :
\[
\lforall[X][(X(x)\liff X(y))].
\]
\end{ex}

\begin{prob}
Montrez que
$\lforall[X][(X(x)\lif X(y))]$
définit $\Id{\Domain{M}}$.
Attention : il s'agit de $\lif$, non de~$\liff$!
\end{prob}

\begin{ex}
Si $R$ est un !!{predicate} binaire,
$\Assign{R}{M}$ est une relation binaire
sur~$\Domain{M}$. Au risque d'une certaine
confusion, nous noterons~$R$ aussi bien le
!!{predicate} que la relation
$\Assign{R}{M}$ elle-même. La
\emph{fermeture transitive}~$R^*$ de~$R$
relie $a$ à~$b$ si, pour certains
$c_1$, \dots, $c_k$, on a
$R(a,c_1)$, $R(c_1,c_2)$, \dots, $R(c_k,b)$.
Le cas $k=0$ est inclus : si $R(a,b)$,
alors $R^*(a,b)$. Ainsi $R\subseteq R^*$.
En fait, $R^*$ est la plus petite relation
transitive qui contient~$R$. Nous pouvons exprimer
au second ordre que $X$ est une relation
transitive contenant~$R$ :
\begin{multline*}
!B_R(X)\ident
\lforall[x][\lforall[y][(R(x,y)\lif X(x,y))]]
\land {}\\
\lforall[x][\lforall[y][\lforall[z]
[((X(x,y)\land X(y,z))\lif X(x,z))]]].
\end{multline*}
La première conjonction dit que $R\subseteq X$;
la seconde, que $X$ est transitive.

Dire que $X$ est la plus petite de ces relations,
c'est dire qu'elle est incluse dans toute relation
transitive contenant~$R$. Nous définissons donc
la fermeture transitive de~$R$ par la
!!{formula}
\[
R^*(X)\ident !B_R(X)\land
\lforall[Y][(!B_R(Y)\lif
\lforall[x][\lforall[y][(X(x,y)\lif Y(x,y))]])].
\]
On a $\Sat{M}{R^*(X)}[s]$ si et seulement si
$s(X)=R^*$. Aucune !!{formula} du premier ordre ne définit uniformément
la fermeture transitive de~$R$ dans toutes les !!{structure}s.
\end{ex}

\begin{editorial}
La notation~$R^*$ désigne ici la plus petite
relation \emph{transitive} contenant~$R$ :
les chemins considérés ont au moins une arête.
Certaines conventions réservent cette notation à
la fermeture réflexive et transitive; ce n'est
pas le sens adopté par la source dans cet exemple.
\end{editorial}

\end{document}
